A \code{RegisterPool} represents a collection of registers.
It can be used to get or set all registers in a Thread at once.
It satisfies the C++
\href{https://en.cppreference.com/cpp/named_req/ForwardIterator}{LegacyForwardIterator}
concept. It also offers the standard \code{begin/end} members for iteration.

\begin{apient}
LibraryPool::iterator find(Dyninst::MachRegister r)
\end{apient}

\apidesc{
This function returns an iterator that points to the element in the register
pool that equals register r. If r is not found, then this function returns
end().
}

\begin{apient}
size_t size() const
\end{apient}

\apidesc{
This function returns the number of elements in the register set.
}

\begin{apient}
Dyninst::MachRegisterVal& operator[](Dyninst::MachRegister r)
\end{apient}

\apidesc{
This function returns a reference to the value associated with the register r in
this register pool. It can be used to efficiently get and set the values of
registers in this pool, or to fill the pool with new
MachRegisters.
}